\FloatBarrier
\section{本章小结}
\label{sec:direct-least-squares-summary}

本章回答了章首的问题：求解目标是最小二乘解或最小范数解时，怎样不求逆、不显式形成不必要的矩阵，在有限步内把选定的解算出来。三条路线的区别在于怎样处理$\bm A$。

正规方程法利用残差正交化为对称正定方程，再用上一章的Cholesky分解与三角代入求解；它可以按行累计数据，但使条件数平方，只适合条件较好的满列秩问题。QR求解用正交变换保持残差长度，把目标分成可消去的部分和不可拟合的部分，满列秩时回代$\bm R\hat{\bm x}=\bm c_1$，满行秩时对$\bm A^\top$分解得到最小范数解。

三种QR实现的差别在于怎样计算正交因子。修正Gram--Schmidt直接算出正交基，有限精度下可能失去正交性；Householder反射一次消去整列，选取同号符号避免相消，是稠密问题的标准做法；Givens旋转一次消去一个元素，适合稀疏结构和逐行更新。秩亏或秩未知时，SVD求解在奇异坐标下把目标分离成独立的标量问题，对任意形状和秩给出$\bm A^\dagger\bm b$；加权问题经白化化为普通情形，齐次问题取完整右奇异基中伸缩量最小的方向。

无论用哪种方法，结果都要回到原始数据上按目标检验：最小二乘解看残差是否与列空间正交，最小范数解还要看是否与零空间正交。这些检验说明算出的$\bm x$符合目标，却不说明它离真解有多近。条件数平方和小奇异值的放大在本章只作了算法层面的说明，结果的可信程度、数值秩的判定以及正则化留给第~\ref{chap:solution-error-regularization}章。矩阵太大而无法分解，或只能通过矩阵与向量的乘积访问矩阵时，就需要从第~\ref{chap:stationary-iterations}章开始的迭代法。
